\documentclass[11pt]{letter}
\usepackage[utf8]{inputenc}
\usepackage[T1]{fontenc}
\usepackage{lmodern}
\usepackage{hyperref}

\signature{Qingyong Wang and Lichuan Gu}
\address{School of Artificial Intelligence\\
Anhui Agricultural University\\
Hefei 230036, China\\
Email: wqy@ahau.edu.cn; glc@ahau.edu.cn}

\begin{document}

\begin{letter}{Editorial Office\\
Bioinformatics\\
Oxford University Press\\
Great Clarendon Street\\
Oxford OX2 6DP\\
United Kingdom}

\opening{Dear Editor,}

We are pleased to submit our manuscript entitled:

\textbf{``Large-scale paired prediction of GPCR--G protein coupling specificity with protein language models and topology-aware features''}

for consideration as an \textbf{Original Paper} in \textit{Bioinformatics}.

\textbf{Summary}

Predicting GPCR--G protein coupling specificity is a long-standing challenge in computational biology, with direct implications for orphan receptor deorphanization and rational drug design. In this study, we present a systematic computational framework that reformulates the problem as paired (GPCR, G protein family) binary classification, using a curated dataset of 1,647 experimentally annotated pairs spanning 431 GPCRs and 4 G protein families.

Our main contributions are:

\begin{enumerate}
    \item \textbf{Systematic architecture comparison:} We benchmark SVM (RBF), cross-attention deep networks, multi-task learning, MLP, Random Forest, and XGBoost under identical conditions, using ESM-2 embeddings at two scales (8M, 650M).

    \item \textbf{Topology-aware ICL features:} Using UniProt transmembrane annotations, we extract ICL2/3 local features and demonstrate a critical \textbf{dimension alignment requirement}: local features must match the global embedding dimension.

    \item \textbf{Negative result with practical impact:} 38-dimensional AlphaFold structural descriptors provide no incremental benefit beyond ESM-2 650M + ICL, suggesting that high-capacity PLMs implicitly encode sufficient structural information---a finding with practical implications for future work.

    \item \textbf{Multi-task analysis:} We show that removing G protein identity information degrades AUC from 0.862 to 0.802, confirming the paired formulation is essential.
\end{enumerate}

The best-performing model (cross-attention + 650M ESM-2 + ICL features) achieves an AUC of 0.8619 $\pm$ 0.0249 under strict cluster-aware cross-validation, significantly outperforming SVM (0.8331).

\textbf{Why Bioinformatics?}

We believe this work is well-suited for \textit{Bioinformatics} because it provides (i) a rigorous benchmarking framework for GPCR--G protein coupling prediction, (ii) generalizable methodological insights (dimension alignment, AlphaFold redundancy) that apply beyond this specific task, and (iii) fully reproducible code and data.

\textbf{Data and Code Availability}

All code, data, and documentation are publicly available at \url{https://github.com/nblvguohao/topology-aware-gpcr-coupling} and archived on Zenodo.

\textbf{Conflict of Interest}

The authors declare no competing interests.

\closing{Thank you for considering our submission. We look forward to your response.}

\vspace{1cm}

Sincerely,

\vspace{0.5cm}

\textbf{Qingyong Wang and Lichuan Gu}

On behalf of all co-authors

\end{letter}

\end{document}
